% Options for packages loaded elsewhere
\PassOptionsToPackage{unicode}{hyperref}
\PassOptionsToPackage{hyphens}{url}
\PassOptionsToPackage{dvipsnames,svgnames,x11names}{xcolor}
%
\documentclass[
  11pt,
  a4paper,
  oneside,
  openany]{scrbook}

\usepackage{amsmath,amssymb}
\usepackage{setspace}
\usepackage{iftex}
\ifPDFTeX
  \usepackage[T1]{fontenc}
  \usepackage[utf8]{inputenc}
  \usepackage{textcomp} % provide euro and other symbols
\else % if luatex or xetex
  \usepackage{unicode-math}
  \defaultfontfeatures{Scale=MatchLowercase}
  \defaultfontfeatures[\rmfamily]{Ligatures=TeX,Scale=1}
\fi
\usepackage{lmodern}
\ifPDFTeX\else  
    % xetex/luatex font selection
\fi
% Use upquote if available, for straight quotes in verbatim environments
\IfFileExists{upquote.sty}{\usepackage{upquote}}{}
\IfFileExists{microtype.sty}{% use microtype if available
  \usepackage[]{microtype}
  \UseMicrotypeSet[protrusion]{basicmath} % disable protrusion for tt fonts
}{}
\makeatletter
\@ifundefined{KOMAClassName}{% if non-KOMA class
  \IfFileExists{parskip.sty}{%
    \usepackage{parskip}
  }{% else
    \setlength{\parindent}{0pt}
    \setlength{\parskip}{6pt plus 2pt minus 1pt}}
}{% if KOMA class
  \KOMAoptions{parskip=half}}
\makeatother
\usepackage{xcolor}
\usepackage[top=25mm,bottom=25mm,left=30mm,right=25mm]{geometry}
\setlength{\emergencystretch}{3em} % prevent overfull lines
\setcounter{secnumdepth}{5}
% Make \paragraph and \subparagraph free-standing
\makeatletter
\ifx\paragraph\undefined\else
  \let\oldparagraph\paragraph
  \renewcommand{\paragraph}{
    \@ifstar
      \xxxParagraphStar
      \xxxParagraphNoStar
  }
  \newcommand{\xxxParagraphStar}[1]{\oldparagraph*{#1}\mbox{}}
  \newcommand{\xxxParagraphNoStar}[1]{\oldparagraph{#1}\mbox{}}
\fi
\ifx\subparagraph\undefined\else
  \let\oldsubparagraph\subparagraph
  \renewcommand{\subparagraph}{
    \@ifstar
      \xxxSubParagraphStar
      \xxxSubParagraphNoStar
  }
  \newcommand{\xxxSubParagraphStar}[1]{\oldsubparagraph*{#1}\mbox{}}
  \newcommand{\xxxSubParagraphNoStar}[1]{\oldsubparagraph{#1}\mbox{}}
\fi
\makeatother

\usepackage{color}
\usepackage{fancyvrb}
\newcommand{\VerbBar}{|}
\newcommand{\VERB}{\Verb[commandchars=\\\{\}]}
\DefineVerbatimEnvironment{Highlighting}{Verbatim}{commandchars=\\\{\}}
% Add ',fontsize=\small' for more characters per line
\usepackage{framed}
\definecolor{shadecolor}{RGB}{248,248,248}
\newenvironment{Shaded}{\begin{snugshade}}{\end{snugshade}}
\newcommand{\AlertTok}[1]{\textcolor[rgb]{0.94,0.16,0.16}{#1}}
\newcommand{\AnnotationTok}[1]{\textcolor[rgb]{0.56,0.35,0.01}{\textbf{\textit{#1}}}}
\newcommand{\AttributeTok}[1]{\textcolor[rgb]{0.13,0.29,0.53}{#1}}
\newcommand{\BaseNTok}[1]{\textcolor[rgb]{0.00,0.00,0.81}{#1}}
\newcommand{\BuiltInTok}[1]{#1}
\newcommand{\CharTok}[1]{\textcolor[rgb]{0.31,0.60,0.02}{#1}}
\newcommand{\CommentTok}[1]{\textcolor[rgb]{0.56,0.35,0.01}{\textit{#1}}}
\newcommand{\CommentVarTok}[1]{\textcolor[rgb]{0.56,0.35,0.01}{\textbf{\textit{#1}}}}
\newcommand{\ConstantTok}[1]{\textcolor[rgb]{0.56,0.35,0.01}{#1}}
\newcommand{\ControlFlowTok}[1]{\textcolor[rgb]{0.13,0.29,0.53}{\textbf{#1}}}
\newcommand{\DataTypeTok}[1]{\textcolor[rgb]{0.13,0.29,0.53}{#1}}
\newcommand{\DecValTok}[1]{\textcolor[rgb]{0.00,0.00,0.81}{#1}}
\newcommand{\DocumentationTok}[1]{\textcolor[rgb]{0.56,0.35,0.01}{\textbf{\textit{#1}}}}
\newcommand{\ErrorTok}[1]{\textcolor[rgb]{0.64,0.00,0.00}{\textbf{#1}}}
\newcommand{\ExtensionTok}[1]{#1}
\newcommand{\FloatTok}[1]{\textcolor[rgb]{0.00,0.00,0.81}{#1}}
\newcommand{\FunctionTok}[1]{\textcolor[rgb]{0.13,0.29,0.53}{\textbf{#1}}}
\newcommand{\ImportTok}[1]{#1}
\newcommand{\InformationTok}[1]{\textcolor[rgb]{0.56,0.35,0.01}{\textbf{\textit{#1}}}}
\newcommand{\KeywordTok}[1]{\textcolor[rgb]{0.13,0.29,0.53}{\textbf{#1}}}
\newcommand{\NormalTok}[1]{#1}
\newcommand{\OperatorTok}[1]{\textcolor[rgb]{0.81,0.36,0.00}{\textbf{#1}}}
\newcommand{\OtherTok}[1]{\textcolor[rgb]{0.56,0.35,0.01}{#1}}
\newcommand{\PreprocessorTok}[1]{\textcolor[rgb]{0.56,0.35,0.01}{\textit{#1}}}
\newcommand{\RegionMarkerTok}[1]{#1}
\newcommand{\SpecialCharTok}[1]{\textcolor[rgb]{0.81,0.36,0.00}{\textbf{#1}}}
\newcommand{\SpecialStringTok}[1]{\textcolor[rgb]{0.31,0.60,0.02}{#1}}
\newcommand{\StringTok}[1]{\textcolor[rgb]{0.31,0.60,0.02}{#1}}
\newcommand{\VariableTok}[1]{\textcolor[rgb]{0.00,0.00,0.00}{#1}}
\newcommand{\VerbatimStringTok}[1]{\textcolor[rgb]{0.31,0.60,0.02}{#1}}
\newcommand{\WarningTok}[1]{\textcolor[rgb]{0.56,0.35,0.01}{\textbf{\textit{#1}}}}

\providecommand{\tightlist}{%
  \setlength{\itemsep}{0pt}\setlength{\parskip}{0pt}}\usepackage{longtable,booktabs,array}
\usepackage{calc} % for calculating minipage widths
% Correct order of tables after \paragraph or \subparagraph
\usepackage{etoolbox}
\makeatletter
\patchcmd\longtable{\par}{\if@noskipsec\mbox{}\fi\par}{}{}
\makeatother
% Allow footnotes in longtable head/foot
\IfFileExists{footnotehyper.sty}{\usepackage{footnotehyper}}{\usepackage{footnote}}
\makesavenoteenv{longtable}
\usepackage{graphicx}
\makeatletter
\def\maxwidth{\ifdim\Gin@nat@width>\linewidth\linewidth\else\Gin@nat@width\fi}
\def\maxheight{\ifdim\Gin@nat@height>\textheight\textheight\else\Gin@nat@height\fi}
\makeatother
% Scale images if necessary, so that they will not overflow the page
% margins by default, and it is still possible to overwrite the defaults
% using explicit options in \includegraphics[width, height, ...]{}
\setkeys{Gin}{width=\maxwidth,height=\maxheight,keepaspectratio}
% Set default figure placement to htbp
\makeatletter
\def\fps@figure{htbp}
\makeatother
% definitions for citeproc citations
\NewDocumentCommand\citeproctext{}{}
\NewDocumentCommand\citeproc{mm}{%
  \begingroup\def\citeproctext{#2}\cite{#1}\endgroup}
\makeatletter
 % allow citations to break across lines
 \let\@cite@ofmt\@firstofone
 % avoid brackets around text for \cite:
 \def\@biblabel#1{}
 \def\@cite#1#2{{#1\if@tempswa , #2\fi}}
\makeatother
\newlength{\cslhangindent}
\setlength{\cslhangindent}{1.5em}
\newlength{\csllabelwidth}
\setlength{\csllabelwidth}{3em}
\newenvironment{CSLReferences}[2] % #1 hanging-indent, #2 entry-spacing
 {\begin{list}{}{%
  \setlength{\itemindent}{0pt}
  \setlength{\leftmargin}{0pt}
  \setlength{\parsep}{0pt}
  % turn on hanging indent if param 1 is 1
  \ifodd #1
   \setlength{\leftmargin}{\cslhangindent}
   \setlength{\itemindent}{-1\cslhangindent}
  \fi
  % set entry spacing
  \setlength{\itemsep}{#2\baselineskip}}}
 {\end{list}}
\usepackage{calc}
\newcommand{\CSLBlock}[1]{\hfill\break\parbox[t]{\linewidth}{\strut\ignorespaces#1\strut}}
\newcommand{\CSLLeftMargin}[1]{\parbox[t]{\csllabelwidth}{\strut#1\strut}}
\newcommand{\CSLRightInline}[1]{\parbox[t]{\linewidth - \csllabelwidth}{\strut#1\strut}}
\newcommand{\CSLIndent}[1]{\hspace{\cslhangindent}#1}

% assets/preamble.tex
% Extra LaTeX for the PDF output of the Stan handbook

% ── Fonts ────────────────────────────────────────────────────────────────────
\usepackage{fontenc}
\usepackage{lmodern}

% ── Math ─────────────────────────────────────────────────────────────────────
\usepackage{amsmath, amssymb, bm}

% ── Stan listings ────────────────────────────────────────────────────────────
\usepackage{listings}
\lstdefinelanguage{Stan}{
  keywords={data, parameters, transformed, model, generated, quantities,
            functions, int, real, vector, matrix, array, lower, upper,
            for, in, if, else, while, return, target},
  keywordstyle=\color[HTML]{2e7d6a}\bfseries,
  comment=[l]{//},
  morecomment=[s]{/*}{*/},
  commentstyle=\color{gray}\itshape,
  stringstyle=\color[HTML]{8B0000},
  basicstyle=\ttfamily\small,
  breaklines=true,
  frame=single,
  rulecolor=\color[HTML]{cccccc},
  backgroundcolor=\color[HTML]{f5f7f5}
}

% ── Callout boxes ─────────────────────────────────────────────────────────────
\usepackage{tcolorbox}
\tcbuselibrary{skins, breakable}

% ── Hyperlinks ────────────────────────────────────────────────────────────────
\usepackage{hyperref}
\hypersetup{
  colorlinks = true,
  linkcolor  = {blue!60!black},
  citecolor  = {green!50!black},
  urlcolor   = {blue!80!black}
}
\makeatletter
\@ifpackageloaded{tcolorbox}{}{\usepackage[skins,breakable]{tcolorbox}}
\@ifpackageloaded{fontawesome5}{}{\usepackage{fontawesome5}}
\definecolor{quarto-callout-color}{HTML}{909090}
\definecolor{quarto-callout-note-color}{HTML}{0758E5}
\definecolor{quarto-callout-important-color}{HTML}{CC1914}
\definecolor{quarto-callout-warning-color}{HTML}{EB9113}
\definecolor{quarto-callout-tip-color}{HTML}{00A047}
\definecolor{quarto-callout-caution-color}{HTML}{FC5300}
\definecolor{quarto-callout-color-frame}{HTML}{acacac}
\definecolor{quarto-callout-note-color-frame}{HTML}{4582ec}
\definecolor{quarto-callout-important-color-frame}{HTML}{d9534f}
\definecolor{quarto-callout-warning-color-frame}{HTML}{f0ad4e}
\definecolor{quarto-callout-tip-color-frame}{HTML}{02b875}
\definecolor{quarto-callout-caution-color-frame}{HTML}{fd7e14}
\makeatother
\makeatletter
\@ifpackageloaded{bookmark}{}{\usepackage{bookmark}}
\makeatother
\makeatletter
\@ifpackageloaded{caption}{}{\usepackage{caption}}
\AtBeginDocument{%
\ifdefined\contentsname
  \renewcommand*\contentsname{Table of contents}
\else
  \newcommand\contentsname{Table of contents}
\fi
\ifdefined\listfigurename
  \renewcommand*\listfigurename{List of Figures}
\else
  \newcommand\listfigurename{List of Figures}
\fi
\ifdefined\listtablename
  \renewcommand*\listtablename{List of Tables}
\else
  \newcommand\listtablename{List of Tables}
\fi
\ifdefined\figurename
  \renewcommand*\figurename{Figure}
\else
  \newcommand\figurename{Figure}
\fi
\ifdefined\tablename
  \renewcommand*\tablename{Table}
\else
  \newcommand\tablename{Table}
\fi
}
\@ifpackageloaded{float}{}{\usepackage{float}}
\floatstyle{ruled}
\@ifundefined{c@chapter}{\newfloat{codelisting}{h}{lop}}{\newfloat{codelisting}{h}{lop}[chapter]}
\floatname{codelisting}{Listing}
\newcommand*\listoflistings{\listof{codelisting}{List of Listings}}
\makeatother
\makeatletter
\makeatother
\makeatletter
\@ifpackageloaded{caption}{}{\usepackage{caption}}
\@ifpackageloaded{subcaption}{}{\usepackage{subcaption}}
\makeatother

\ifLuaTeX
  \usepackage{selnolig}  % disable illegal ligatures
\fi
\usepackage{bookmark}

\IfFileExists{xurl.sty}{\usepackage{xurl}}{} % add URL line breaks if available
\urlstyle{same} % disable monospaced font for URLs
\hypersetup{
  pdftitle={Writing Statistical Models in Stan},
  pdfauthor={Anicet Ebou},
  colorlinks=true,
  linkcolor={blue!60!black},
  filecolor={Maroon},
  citecolor={Blue},
  urlcolor={Blue},
  pdfcreator={LaTeX via pandoc}}


\title{Writing Statistical Models in Stan}
\usepackage{etoolbox}
\makeatletter
\providecommand{\subtitle}[1]{% add subtitle to \maketitle
  \apptocmd{\@title}{\par {\large #1 \par}}{}{}
}
\makeatother
\subtitle{A Practical Handbook for Life Scientists}
\author{Anicet Ebou}
\date{2026-06-16}

\begin{document}
\frontmatter
\maketitle

\renewcommand*\contentsname{Table of contents}
{
\hypersetup{linkcolor=}
\setcounter{tocdepth}{1}
\tableofcontents
}

\setstretch{1.25}
\mainmatter
\bookmarksetup{startatroot}

\chapter*{Preface}\label{preface}
\addcontentsline{toc}{chapter}{Preface}

\markboth{Preface}{Preface}

This handbook teaches you to write statistical models directly in
\textbf{Stan} --- the probabilistic programming language behind many
modern Bayesian workflows.

\section*{Who this book is for}\label{who-this-book-is-for}
\addcontentsline{toc}{section}{Who this book is for}

\markright{Who this book is for}

You are a life scientist, agronomist, or ecologist who:

\begin{itemize}
\tightlist
\item
  is comfortable fitting linear models and ANOVA in R,
\item
  has heard of Bayesian statistics but found it abstract,
\item
  wants to move beyond black-box software and understand \emph{what the
  model is actually doing}.
\end{itemize}

No prior knowledge of Stan or Bayesian inference is assumed. Familiarity
with R and \texttt{lm()} / \texttt{aov()} is sufficient to begin.

\section*{How to use this handbook}\label{how-to-use-this-handbook}
\addcontentsline{toc}{section}{How to use this handbook}

\markright{How to use this handbook}

Each chapter introduces one model family, shows a complete Stan program,
and walks through interpretation and diagnostics. \textbf{Every code
block is runnable} --- copy it into your own script and experiment.

The chapters in \textbf{Part II} (linear models and ANOVA) are the heart
of the book. Parts III--IV build on them progressively. Part V gives you
the habits that separate a competent Stan user from a self-sufficient
one.

\section*{Conventions}\label{conventions}
\addcontentsline{toc}{section}{Conventions}

\markright{Conventions}

\begin{longtable}[]{@{}
  >{\raggedright\arraybackslash}p{(\columnwidth - 2\tabcolsep) * \real{0.3000}}
  >{\raggedright\arraybackslash}p{(\columnwidth - 2\tabcolsep) * \real{0.7000}}@{}}
\caption{Typographic conventions used in this book.}\tabularnewline
\toprule\noalign{}
\begin{minipage}[b]{\linewidth}\raggedright
Convention
\end{minipage} & \begin{minipage}[b]{\linewidth}\raggedright
Meaning
\end{minipage} \\
\midrule\noalign{}
\endfirsthead
\toprule\noalign{}
\begin{minipage}[b]{\linewidth}\raggedright
Convention
\end{minipage} & \begin{minipage}[b]{\linewidth}\raggedright
Meaning
\end{minipage} \\
\midrule\noalign{}
\endhead
\bottomrule\noalign{}
\endlastfoot
\texttt{code\ font} & R or Stan code, file names, function names \\
\emph{italics} & mathematical quantities or new terms at first use \\
\textbf{bold} & key concepts \\
::: \{.callout-tip\} blocks & practical advice worth remembering \\
::: \{.callout-warning\} blocks & common mistakes to avoid \\
\end{longtable}

Mathematical notation follows standard convention: scalars in lower-case
italics (\(y\), \(\mu\), \(\sigma\)), vectors in bold
(\(\boldsymbol{x}\)), matrices in bold upper-case (\(\mathbf{X}\)).

\section*{Software versions}\label{software-versions}
\addcontentsline{toc}{section}{Software versions}

\markright{Software versions}

\begin{longtable}[]{@{}lll@{}}
\caption{Key R packages used throughout this handbook.}\tabularnewline
\toprule\noalign{}
& Package & Version \\
\midrule\noalign{}
\endfirsthead
\toprule\noalign{}
& Package & Version \\
\midrule\noalign{}
\endhead
\bottomrule\noalign{}
\endlastfoot
cmdstanr & cmdstanr & 0.9.0 \\
posterior & posterior & 1.7.0 \\
bayesplot & bayesplot & 1.15.0 \\
loo & loo & 2.9.0 \\
\end{longtable}

Stan version and CmdStan path are reported by
\texttt{cmdstanr::cmdstan\_version()}.

\section*{Acknowledgements}\label{acknowledgements}
\addcontentsline{toc}{section}{Acknowledgements}

\markright{Acknowledgements}

\emph{{[}To be completed.{]}}

\part{Part I - Getting Started}

\chapter{Why Stan?}\label{why-stan}

\section{Overview}\label{overview}

Before writing a single line of Stan code, it is worth asking an honest
question: \emph{why bother?} R already provides \texttt{lm()},
\texttt{aov()}, \texttt{glm()}, and a rich ecosystem of mixed-model
packages. Bayesian wrappers like \textbf{brms} can fit hierarchical
models with a formula interface nearly as simple as \texttt{lmer()}. So
what does writing Stan directly buy you?

This chapter answers that question. It positions Stan in the landscape
of statistical software, explains what probabilistic programming
actually means, and makes the case that the effort of learning Stan pays
dividends that no wrapper can fully replicate.

\section{Learning objectives}\label{learning-objectives}

By the end of this chapter you will be able to:

\begin{itemize}
\tightlist
\item
  Explain the difference between \emph{fitting a model} and
  \emph{specifying a model}.
\item
  Describe what Stan does and how it differs from classical statistical
  software and from higher-level Bayesian wrappers.
\item
  Articulate at least three concrete reasons why writing Stan directly
  is worth the investment for a practicing life scientist.
\end{itemize}

\begin{center}\rule{0.5\linewidth}{0.5pt}\end{center}

\section{The gap between fitting and
specifying}\label{the-gap-between-fitting-and-specifying}

When you call
\texttt{lm(yield\ \textasciitilde{}\ treatment\ +\ block,\ data\ =\ field\_trial)},
you are \emph{fitting} a model --- R finds parameter estimates and
standard errors, and returns them. What you are not doing is
\emph{specifying} a model. The statistical assumptions (normal errors,
constant variance, independence) are baked in. The algorithm (ordinary
least squares) is fixed. There is no mechanism for expressing what you
know about crop yields before seeing the data.

This works well most of the time. For a balanced, replicated field trial
with continuous outcomes, \texttt{lm()} is hard to beat. But
life-science data are rarely this clean:

\begin{itemize}
\tightlist
\item
  Outcomes are often counts, proportions, or presence/absence.
\item
  Observations are grouped (plots within farms, samples within subjects,
  leaves within trees) and genuinely not independent.
\item
  Sample sizes are small and estimates are uncertain in ways that a
  single point estimate and standard error cannot capture.
\item
  Prior knowledge --- about plausible effect sizes, about measurement
  error, about the biology --- exists and should inform the analysis.
\end{itemize}

As soon as one of these complications arises, black-box software either
forces you into a pre-built model family or invites you to ignore the
problem entirely. Stan takes a different stance: \emph{you write the
model, Stan fits it.}

\begin{center}\rule{0.5\linewidth}{0.5pt}\end{center}

\section{What Stan is}\label{what-stan-is}

Stan is a \textbf{probabilistic programming language} designed for
Bayesian statistical modelling (Carpenter et al. 2017). You write a
program that encodes:

\begin{enumerate}
\def\labelenumi{\arabic{enumi}.}
\tightlist
\item
  \textbf{The data} you will observe.
\item
  \textbf{The parameters} you want to learn.
\item
  \textbf{The model} --- how the data arise from the parameters (the
  likelihood) and what you believe about the parameters before seeing
  the data (the prior).
\end{enumerate}

Stan then uses \textbf{Hamiltonian Monte Carlo} (HMC) --- specifically
the No-U-Turn Sampler (NUTS) --- to draw samples from the
\emph{posterior distribution} of your parameters given your data
(Betancourt 2017). Those samples are your inferential output: you can
compute means, medians, credible intervals, predictions, and model
comparisons directly from them.

A Stan program for a simple normal model looks like this:

\begin{Shaded}
\begin{Highlighting}[]
\KeywordTok{data}\NormalTok{ \{}
  \DataTypeTok{int}\NormalTok{\textless{}}\KeywordTok{lower}\NormalTok{=}\DecValTok{1}\NormalTok{\textgreater{} N;       }\CommentTok{// number of observations}
  \DataTypeTok{vector}\NormalTok{[N] y;          }\CommentTok{// outcome variable}
\NormalTok{\}}

\KeywordTok{parameters}\NormalTok{ \{}
  \DataTypeTok{real}\NormalTok{ mu;              }\CommentTok{// population mean}
  \DataTypeTok{real}\NormalTok{\textless{}}\KeywordTok{lower}\NormalTok{=}\DecValTok{0}\NormalTok{\textgreater{} sigma;  }\CommentTok{// standard deviation (must be positive)}
\NormalTok{\}}

\KeywordTok{model}\NormalTok{ \{}
  \CommentTok{// Priors}
\NormalTok{  mu    \textasciitilde{} normal(}\DecValTok{0}\NormalTok{, }\DecValTok{10}\NormalTok{);}
\NormalTok{  sigma \textasciitilde{} exponential(}\DecValTok{1}\NormalTok{);}

  \CommentTok{// Likelihood}
\NormalTok{  y \textasciitilde{} normal(mu, sigma);}
\NormalTok{\}}
\end{Highlighting}
\end{Shaded}

Seven lines of model code. Every assumption is explicit. Every prior is
a deliberate choice, not a hidden default. This is the core promise of
Stan.

\begin{center}\rule{0.5\linewidth}{0.5pt}\end{center}

\section{Stan in the landscape}\label{stan-in-the-landscape}

It helps to know where Stan sits relative to other tools you may have
heard of (Table~\ref{tbl-landscape}).

\begin{longtable}[]{@{}
  >{\raggedright\arraybackslash}p{(\columnwidth - 6\tabcolsep) * \real{0.2500}}
  >{\raggedright\arraybackslash}p{(\columnwidth - 6\tabcolsep) * \real{0.2500}}
  >{\raggedright\arraybackslash}p{(\columnwidth - 6\tabcolsep) * \real{0.2500}}
  >{\raggedright\arraybackslash}p{(\columnwidth - 6\tabcolsep) * \real{0.2500}}@{}}
\caption{Statistical modelling tools and their
trade-offs.}\label{tbl-landscape}\tabularnewline
\toprule\noalign{}
\begin{minipage}[b]{\linewidth}\raggedright
Tool
\end{minipage} & \begin{minipage}[b]{\linewidth}\raggedright
Paradigm
\end{minipage} & \begin{minipage}[b]{\linewidth}\raggedright
Strength
\end{minipage} & \begin{minipage}[b]{\linewidth}\raggedright
Limitation
\end{minipage} \\
\midrule\noalign{}
\endfirsthead
\toprule\noalign{}
\begin{minipage}[b]{\linewidth}\raggedright
Tool
\end{minipage} & \begin{minipage}[b]{\linewidth}\raggedright
Paradigm
\end{minipage} & \begin{minipage}[b]{\linewidth}\raggedright
Strength
\end{minipage} & \begin{minipage}[b]{\linewidth}\raggedright
Limitation
\end{minipage} \\
\midrule\noalign{}
\endhead
\bottomrule\noalign{}
\endlastfoot
\texttt{lm()}, \texttt{aov()} & Frequentist & Fast, exact,
well-understood & Fixed model family, no priors \\
\texttt{lme4}, \texttt{nlme} & Frequentist / REML & Mixed models, widely
used & Convergence issues, limited family \\
JAGS / BUGS & Bayesian (MCMC) & Flexible, mature & Slow, Gibbs sampling,
old syntax \\
\textbf{brms} & Bayesian (Stan backend) & Formula interface, powerful &
Limited when you need a custom model \\
\textbf{Stan} & Bayesian (HMC) & Full flexibility, fast sampler &
Steeper learning curve \\
PyMC & Bayesian (various) & Python ecosystem & Python, not R \\
\end{longtable}

\textbf{brms} deserves special mention because it uses Stan under the
hood. For standard model families --- linear, generalised linear,
multilevel --- brms is excellent and you should use it. The case for
writing Stan directly arises when your model goes beyond what brms's
formula language can express:

\begin{itemize}
\tightlist
\item
  Custom likelihood functions (e.g., zero-inflated models with
  non-standard zero-generating processes).
\item
  Mechanistic models where parameters have biological meaning and
  constraints (e.g., compartmental models, growth curves).
\item
  Spatial models with structured priors (ICAR, Gaussian processes).
\item
  Multi-response models with complex covariance structure.
\item
  Anything where the \emph{parameterisation} matters for sampling
  efficiency.
\end{itemize}

\begin{tcolorbox}[enhanced jigsaw, colframe=quarto-callout-note-color-frame, breakable, title=\textcolor{quarto-callout-note-color}{\faInfo}\hspace{0.5em}{Note}, coltitle=black, bottomtitle=1mm, colbacktitle=quarto-callout-note-color!10!white, toprule=.15mm, toptitle=1mm, leftrule=.75mm, arc=.35mm, titlerule=0mm, rightrule=.15mm, bottomrule=.15mm, opacitybacktitle=0.6, left=2mm, opacityback=0, colback=white]

\textbf{brms and Stan are not competitors.} Many practitioners use brms
for routine models and write Stan directly when the problem demands it.
Learning Stan makes you a better brms user too, because you can read the
Stan code brms generates (via \texttt{brms::stancode()}) and understand
exactly what is being fitted.

\end{tcolorbox}

\begin{center}\rule{0.5\linewidth}{0.5pt}\end{center}

\section{Three reasons to learn Stan}\label{three-reasons-to-learn-stan}

\subsection{1. Your model is your
analysis}\label{your-model-is-your-analysis}

In classical software, the model lives in the function call. In Stan,
the model lives in a file you can read, edit, version-control, and
share. This has practical consequences:

\begin{itemize}
\tightlist
\item
  \textbf{Reviewers and collaborators} can check your statistical
  assumptions directly, not just the software you called.
\item
  \textbf{You} can iterate on the model --- tightening a prior, adding a
  covariate, changing the likelihood --- without hunting for a new
  function.
\item
  \textbf{Reproducibility} is structural: the \texttt{.stan} file is the
  model.
\end{itemize}

\subsection{2. Uncertainty is
first-class}\label{uncertainty-is-first-class}

Stan returns \emph{distributions}, not point estimates. For a slope
parameter \(\beta\), you do not get a single number and a \(p\)-value;
you get thousands of plausible values of \(\beta\) drawn from the
posterior. You can then ask questions that \(p\)-values cannot answer:

\begin{itemize}
\tightlist
\item
  What is the probability that \(\beta > 0\)?
\item
  What is the 90\% credible interval for the expected yield difference
  between treatments?
\item
  If I had to act on this result, what range of outcomes should I plan
  for?
\end{itemize}

This matters especially in applied ecology and agronomy, where the cost
of a wrong decision is real and decisions are made under genuine
uncertainty.

\subsection{3. Small samples, rare events, partial
information}\label{small-samples-rare-events-partial-information}

Bayesian methods --- and Stan in particular --- handle small samples
gracefully because prior information contributes to inference when data
are sparse. A frequentist mixed model may fail to converge with five
observations per group; a Stan model with a sensible prior will return a
posterior that honestly reflects both the data and the prior, with
appropriately wide uncertainty.

This is not a trick or a way to manufacture certainty from thin air. It
is a principled mechanism for combining what you knew before the
experiment with what you learned from it.

\begin{center}\rule{0.5\linewidth}{0.5pt}\end{center}

\section{What this handbook will teach
you}\label{what-this-handbook-will-teach-you}

This handbook takes you from the \texttt{.stan} file for a simple linear
regression all the way to hierarchical models, generalised linear
models, and the practical skills (diagnostics, model comparison,
debugging) that let you use Stan confidently in your own research.

The path follows your existing statistical knowledge:

\begin{verbatim}
lm()  →  Chapter 4 (simple linear regression in Stan)
aov() →  Chapters 6–7 (one-way and two-way ANOVA in Stan)
lmer() → Chapters 8–10 (hierarchical models in Stan)
glm()  → Chapters 11–12 (GLMs in Stan)
\end{verbatim}

Each step translates something you already know into Stan's language. By
the time you reach the end, you will not just know \emph{how} to use
Stan --- you will understand \emph{why} your model works, and you will
be equipped to build models that no wrapper package can express.

\begin{center}\rule{0.5\linewidth}{0.5pt}\end{center}

\section{Summary}\label{summary}

\begin{itemize}
\tightlist
\item
  \textbf{Stan is a probabilistic programming language} for Bayesian
  statistical modelling. You write the model; Stan fits it.
\item
  Unlike classical software, Stan does not fix the model family or the
  algorithm. Every assumption is explicit and every prior is a
  deliberate choice.
\item
  Stan uses \textbf{Hamiltonian Monte Carlo} to draw samples from the
  posterior distribution of your parameters.
\item
  \textbf{brms} covers most standard models via a formula interface.
  Stan is the right tool when your model goes beyond what brms can
  express, or when you need full control over parameterisation.
\item
  Three core benefits: (1) the model is a legible artefact; (2)
  uncertainty is first-class; (3) principled handling of small samples
  and prior information.
\end{itemize}

\begin{center}\rule{0.5\linewidth}{0.5pt}\end{center}

\section{Exercises}\label{exercises}

\begin{enumerate}
\def\labelenumi{\arabic{enumi}.}
\item
  Open the help page for \texttt{brms::stancode()}. Fit a simple linear
  regression with
  \texttt{brm(y\ \textasciitilde{}\ x,\ data\ =\ mydata,\ family\ =\ gaussian())}
  and then call \texttt{stancode()} on the result. Read the generated
  Stan program. Which blocks do you recognise? Which parts are
  unfamiliar?
\item
  Think of a modelling problem from your own field where \texttt{lm()}
  or \texttt{aov()} feels inadequate --- perhaps because of grouped
  data, a non-normal outcome, or small sample sizes. Write a short
  paragraph describing the problem. We will return to problems like this
  throughout the handbook.
\item
  Stan's website (\url{https://mc-stan.org}) hosts a \emph{Case Studies}
  section. Browse the titles and find one that relates to ecology,
  agriculture, or biology. Read the introduction and note: (a) what the
  model is trying to infer, and (b) why a standard approach would be
  insufficient.
\end{enumerate}

\begin{center}\rule{0.5\linewidth}{0.5pt}\end{center}

\section{Further reading}\label{further-reading}

\begin{itemize}
\tightlist
\item
  McElreath (2020, Ch. 1) --- a compelling philosophical case for
  Bayesian modelling, aimed squarely at life scientists.
\item
  Gelman, Hill, and Vehtari (2020, Ch. 1--2) --- contextualises
  regression as a general inferential tool, not just a prediction
  machine.
\item
  Betancourt (2017) --- a conceptual introduction to HMC for readers who
  want to understand \emph{why} Stan's sampler is better than Gibbs
  sampling.
\item
  Stan documentation: \url{https://mc-stan.org/docs/stan-users-guide/}
\end{itemize}

\chapter{Setting Up Your Environment}\label{setting-up-your-environment}

\section{Overview}\label{overview-1}

Installing CmdStanR, verifying the toolchain, and running your first
Stan program.

\section{Learning objectives}\label{learning-objectives-1}

By the end of this chapter you will be able to:

\begin{itemize}
\tightlist
\item[$\square$]
  Objective 1
\item[$\square$]
  Objective 2
\item[$\square$]
  Objective 3
\end{itemize}

\section{Content}\label{content}

\emph{Content to be written.}

\section{Exercises}\label{exercises-1}

\begin{enumerate}
\def\labelenumi{\arabic{enumi}.}
\tightlist
\item
  Exercise 1
\item
  Exercise 2
\end{enumerate}

\section{Summary}\label{summary-1}

\emph{Summary to be written.}

\chapter{Thinking Like a Bayesian}\label{thinking-like-a-bayesian}

\section{Overview}\label{overview-2}

Likelihood, prior, posterior, and MCMC intuition --- without the measure
theory.

\section{Learning objectives}\label{learning-objectives-2}

By the end of this chapter you will be able to:

\begin{itemize}
\tightlist
\item[$\square$]
  Objective 1
\item[$\square$]
  Objective 2
\item[$\square$]
  Objective 3
\end{itemize}

\section{Content}\label{content-1}

\emph{Content to be written.}

\section{Exercises}\label{exercises-2}

\begin{enumerate}
\def\labelenumi{\arabic{enumi}.}
\tightlist
\item
  Exercise 1
\item
  Exercise 2
\end{enumerate}

\section{Summary}\label{summary-2}

\emph{Summary to be written.}

\part{Part II - The Linear World}

\chapter{Simple Linear Regression in
Stan}\label{simple-linear-regression-in-stan}

\section{Overview}\label{overview-3}

Writing the data, parameters, and model blocks for a basic regression.

\section{Learning objectives}\label{learning-objectives-3}

By the end of this chapter you will be able to:

\begin{itemize}
\tightlist
\item[$\square$]
  Objective 1
\item[$\square$]
  Objective 2
\item[$\square$]
  Objective 3
\end{itemize}

\section{Content}\label{content-2}

\emph{Content to be written.}

\section{Exercises}\label{exercises-3}

\begin{enumerate}
\def\labelenumi{\arabic{enumi}.}
\tightlist
\item
  Exercise 1
\item
  Exercise 2
\end{enumerate}

\section{Summary}\label{summary-3}

\emph{Summary to be written.}

\chapter{Multiple Regression and
Predictors}\label{multiple-regression-and-predictors}

\section{Overview}\label{overview-4}

Centring predictors, vectorised notation, and posterior predictive
checks.

\section{Learning objectives}\label{learning-objectives-4}

By the end of this chapter you will be able to:

\begin{itemize}
\tightlist
\item[$\square$]
  Objective 1
\item[$\square$]
  Objective 2
\item[$\square$]
  Objective 3
\end{itemize}

\section{Content}\label{content-3}

\emph{Content to be written.}

\section{Exercises}\label{exercises-4}

\begin{enumerate}
\def\labelenumi{\arabic{enumi}.}
\tightlist
\item
  Exercise 1
\item
  Exercise 2
\end{enumerate}

\section{Summary}\label{summary-4}

\emph{Summary to be written.}

\chapter{One-Way ANOVA in Stan}\label{one-way-anova-in-stan}

\section{Overview}\label{overview-5}

Treatment coding, sum-to-zero constraints, and reading ANOVA tables from
Stan output.

\section{Learning objectives}\label{learning-objectives-5}

By the end of this chapter you will be able to:

\begin{itemize}
\tightlist
\item[$\square$]
  Objective 1
\item[$\square$]
  Objective 2
\item[$\square$]
  Objective 3
\end{itemize}

\section{Content}\label{content-4}

\emph{Content to be written.}

\section{Exercises}\label{exercises-5}

\begin{enumerate}
\def\labelenumi{\arabic{enumi}.}
\tightlist
\item
  Exercise 1
\item
  Exercise 2
\end{enumerate}

\section{Summary}\label{summary-5}

\emph{Summary to be written.}

\chapter{Two-Way ANOVA and
Interactions}\label{two-way-anova-and-interactions}

\section{Overview}\label{overview-6}

Additive vs interaction models, identifiability, and visualising
interaction posteriors.

\section{Learning objectives}\label{learning-objectives-6}

By the end of this chapter you will be able to:

\begin{itemize}
\tightlist
\item[$\square$]
  Objective 1
\item[$\square$]
  Objective 2
\item[$\square$]
  Objective 3
\end{itemize}

\section{Content}\label{content-5}

\emph{Content to be written.}

\section{Exercises}\label{exercises-6}

\begin{enumerate}
\def\labelenumi{\arabic{enumi}.}
\tightlist
\item
  Exercise 1
\item
  Exercise 2
\end{enumerate}

\section{Summary}\label{summary-6}

\emph{Summary to be written.}

\part{Part III - Hierarchical Models}

\chapter{Introduction to Hierarchical
Models}\label{introduction-to-hierarchical-models}

\section{Overview}\label{overview-7}

Grouped data, partial pooling, and the random-intercept model.

\section{Learning objectives}\label{learning-objectives-7}

By the end of this chapter you will be able to:

\begin{itemize}
\tightlist
\item[$\square$]
  Objective 1
\item[$\square$]
  Objective 2
\item[$\square$]
  Objective 3
\end{itemize}

\section{Content}\label{content-6}

\emph{Content to be written.}

\section{Exercises}\label{exercises-7}

\begin{enumerate}
\def\labelenumi{\arabic{enumi}.}
\tightlist
\item
  Exercise 1
\item
  Exercise 2
\end{enumerate}

\section{Summary}\label{summary-7}

\emph{Summary to be written.}

\chapter{Random Slopes and Varying
Effects}\label{random-slopes-and-varying-effects}

\section{Overview}\label{overview-8}

Multivariate normal priors, Cholesky parameterisation, and
repeated-measures examples.

\section{Learning objectives}\label{learning-objectives-8}

By the end of this chapter you will be able to:

\begin{itemize}
\tightlist
\item[$\square$]
  Objective 1
\item[$\square$]
  Objective 2
\item[$\square$]
  Objective 3
\end{itemize}

\section{Content}\label{content-7}

\emph{Content to be written.}

\section{Exercises}\label{exercises-8}

\begin{enumerate}
\def\labelenumi{\arabic{enumi}.}
\tightlist
\item
  Exercise 1
\item
  Exercise 2
\end{enumerate}

\section{Summary}\label{summary-8}

\emph{Summary to be written.}

\chapter{Non-Centred
Parameterisation}\label{non-centred-parameterisation}

\section{Overview}\label{overview-9}

The funnel problem, divergent transitions, and how to reparameterise for
better sampling.

\section{Learning objectives}\label{learning-objectives-9}

By the end of this chapter you will be able to:

\begin{itemize}
\tightlist
\item[$\square$]
  Objective 1
\item[$\square$]
  Objective 2
\item[$\square$]
  Objective 3
\end{itemize}

\section{Content}\label{content-8}

\emph{Content to be written.}

\section{Exercises}\label{exercises-9}

\begin{enumerate}
\def\labelenumi{\arabic{enumi}.}
\tightlist
\item
  Exercise 1
\item
  Exercise 2
\end{enumerate}

\section{Summary}\label{summary-9}

\emph{Summary to be written.}

\part{Part IV - Beyond the Normal}

\chapter{Generalised Linear Models in
Stan}\label{generalised-linear-models-in-stan}

\section{Overview}\label{overview-10}

Logistic, Poisson, and negative binomial models with generated
quantities.

\section{Learning objectives}\label{learning-objectives-10}

By the end of this chapter you will be able to:

\begin{itemize}
\tightlist
\item[$\square$]
  Objective 1
\item[$\square$]
  Objective 2
\item[$\square$]
  Objective 3
\end{itemize}

\section{Content}\label{content-9}

\emph{Content to be written.}

\section{Exercises}\label{exercises-10}

\begin{enumerate}
\def\labelenumi{\arabic{enumi}.}
\tightlist
\item
  Exercise 1
\item
  Exercise 2
\end{enumerate}

\section{Summary}\label{summary-10}

\emph{Summary to be written.}

\chapter{Ordinal and Categorical
Outcomes}\label{ordinal-and-categorical-outcomes}

\section{Overview}\label{overview-11}

Ordered logistic regression and the softmax / categorical likelihood.

\section{Learning objectives}\label{learning-objectives-11}

By the end of this chapter you will be able to:

\begin{itemize}
\tightlist
\item[$\square$]
  Objective 1
\item[$\square$]
  Objective 2
\item[$\square$]
  Objective 3
\end{itemize}

\section{Content}\label{content-10}

\emph{Content to be written.}

\section{Exercises}\label{exercises-11}

\begin{enumerate}
\def\labelenumi{\arabic{enumi}.}
\tightlist
\item
  Exercise 1
\item
  Exercise 2
\end{enumerate}

\section{Summary}\label{summary-11}

\emph{Summary to be written.}

\part{Part V - Good Practice and Going Further}

\chapter{Prior Predictive Checks}\label{prior-predictive-checks}

\section{Overview}\label{overview-12}

Simulating from the prior before fitting to catch unreasonable
assumptions early.

\section{Learning objectives}\label{learning-objectives-12}

By the end of this chapter you will be able to:

\begin{itemize}
\tightlist
\item[$\square$]
  Objective 1
\item[$\square$]
  Objective 2
\item[$\square$]
  Objective 3
\end{itemize}

\section{Content}\label{content-11}

\emph{Content to be written.}

\section{Exercises}\label{exercises-12}

\begin{enumerate}
\def\labelenumi{\arabic{enumi}.}
\tightlist
\item
  Exercise 1
\item
  Exercise 2
\end{enumerate}

\section{Summary}\label{summary-12}

\emph{Summary to be written.}

\chapter{Model Comparison and
Selection}\label{model-comparison-and-selection}

\section{Overview}\label{overview-13}

LOO-CV with the loo package, WAIC, and when not to compare models.

\section{Learning objectives}\label{learning-objectives-13}

By the end of this chapter you will be able to:

\begin{itemize}
\tightlist
\item[$\square$]
  Objective 1
\item[$\square$]
  Objective 2
\item[$\square$]
  Objective 3
\end{itemize}

\section{Content}\label{content-12}

\emph{Content to be written.}

\section{Exercises}\label{exercises-13}

\begin{enumerate}
\def\labelenumi{\arabic{enumi}.}
\tightlist
\item
  Exercise 1
\item
  Exercise 2
\end{enumerate}

\section{Summary}\label{summary-13}

\emph{Summary to be written.}

\chapter{Debugging Stan Models}\label{debugging-stan-models}

\section{Overview}\label{overview-14}

Common error messages, divergent transitions, and unit-testing Stan
functions.

\section{Learning objectives}\label{learning-objectives-14}

By the end of this chapter you will be able to:

\begin{itemize}
\tightlist
\item[$\square$]
  Objective 1
\item[$\square$]
  Objective 2
\item[$\square$]
  Objective 3
\end{itemize}

\section{Content}\label{content-13}

\emph{Content to be written.}

\section{Exercises}\label{exercises-14}

\begin{enumerate}
\def\labelenumi{\arabic{enumi}.}
\tightlist
\item
  Exercise 1
\item
  Exercise 2
\end{enumerate}

\section{Summary}\label{summary-14}

\emph{Summary to be written.}

\chapter{From Stan to Publication}\label{from-stan-to-publication}

\section{Overview}\label{overview-15}

Reporting Bayesian results, reproducibility practices, and sharing Stan
code.

\section{Learning objectives}\label{learning-objectives-15}

By the end of this chapter you will be able to:

\begin{itemize}
\tightlist
\item[$\square$]
  Objective 1
\item[$\square$]
  Objective 2
\item[$\square$]
  Objective 3
\end{itemize}

\section{Content}\label{content-14}

\emph{Content to be written.}

\section{Exercises}\label{exercises-15}

\begin{enumerate}
\def\labelenumi{\arabic{enumi}.}
\tightlist
\item
  Exercise 1
\item
  Exercise 2
\end{enumerate}

\section{Summary}\label{summary-15}

\emph{Summary to be written.}

\part{Appendices}

\chapter{Stan Language Quick
Reference}\label{stan-language-quick-reference}

\section{Overview}\label{overview-16}

Block structure, data types, distributions, and vectorisation cheat
sheet.

\section{Learning objectives}\label{learning-objectives-16}

By the end of this chapter you will be able to:

\begin{itemize}
\tightlist
\item[$\square$]
  Objective 1
\item[$\square$]
  Objective 2
\item[$\square$]
  Objective 3
\end{itemize}

\section{Content}\label{content-15}

\emph{Content to be written.}

\section{Exercises}\label{exercises-16}

\begin{enumerate}
\def\labelenumi{\arabic{enumi}.}
\tightlist
\item
  Exercise 1
\item
  Exercise 2
\end{enumerate}

\section{Summary}\label{summary-16}

\emph{Summary to be written.}

\chapter{Prior Distribution Gallery}\label{prior-distribution-gallery}

\section{Overview}\label{overview-17}

Visualisations and guidance for commonly used prior families.

\section{Learning objectives}\label{learning-objectives-17}

By the end of this chapter you will be able to:

\begin{itemize}
\tightlist
\item[$\square$]
  Objective 1
\item[$\square$]
  Objective 2
\item[$\square$]
  Objective 3
\end{itemize}

\section{Content}\label{content-16}

\emph{Content to be written.}

\section{Exercises}\label{exercises-17}

\begin{enumerate}
\def\labelenumi{\arabic{enumi}.}
\tightlist
\item
  Exercise 1
\item
  Exercise 2
\end{enumerate}

\section{Summary}\label{summary-17}

\emph{Summary to be written.}

\chapter{Glossary of MCMC
Diagnostics}\label{glossary-of-mcmc-diagnostics}

\section{Overview}\label{overview-18}

R-hat, ESS, BFMI, divergences, and energy plots explained.

\section{Learning objectives}\label{learning-objectives-18}

By the end of this chapter you will be able to:

\begin{itemize}
\tightlist
\item[$\square$]
  Objective 1
\item[$\square$]
  Objective 2
\item[$\square$]
  Objective 3
\end{itemize}

\section{Content}\label{content-17}

\emph{Content to be written.}

\section{Exercises}\label{exercises-18}

\begin{enumerate}
\def\labelenumi{\arabic{enumi}.}
\tightlist
\item
  Exercise 1
\item
  Exercise 2
\end{enumerate}

\section{Summary}\label{summary-18}

\emph{Summary to be written.}

\chapter{Further Reading and
Resources}\label{further-reading-and-resources}

\section{Overview}\label{overview-19}

Books, papers, Stan forums, and online courses to deepen your
understanding.

\section{Learning objectives}\label{learning-objectives-19}

By the end of this chapter you will be able to:

\begin{itemize}
\tightlist
\item[$\square$]
  Objective 1
\item[$\square$]
  Objective 2
\item[$\square$]
  Objective 3
\end{itemize}

\section{Content}\label{content-18}

\emph{Content to be written.}

\section{Exercises}\label{exercises-19}

\begin{enumerate}
\def\labelenumi{\arabic{enumi}.}
\tightlist
\item
  Exercise 1
\item
  Exercise 2
\end{enumerate}

\section{Summary}\label{summary-19}

\emph{Summary to be written.}

\cleardoublepage
\phantomsection
\addcontentsline{toc}{part}{Appendices}
\appendix

\chapter*{References}\label{references}
\addcontentsline{toc}{chapter}{References}

\markboth{References}{References}

\phantomsection\label{refs}
\begin{CSLReferences}{1}{0}
\bibitem[\citeproctext]{ref-betancourt2017conceptual}
Betancourt, Michael. 2017. {``A Conceptual Introduction to {Hamiltonian
Monte Carlo}.''} \emph{arXiv Preprint arXiv:1701.02434}.

\bibitem[\citeproctext]{ref-stan2017}
Carpenter, Bob, Andrew Gelman, Matthew D Hoffman, Daniel Lee, Ben
Goodrich, Michael Betancourt, Marcus Brubaker, Jiqiang Guo, Peter Li,
and Allen Riddell. 2017. {``Stan: A Probabilistic Programming
Language.''} \emph{Journal of Statistical Software} 76 (1): 1--32.
\url{https://doi.org/10.18637/jss.v076.i01}.

\bibitem[\citeproctext]{ref-gelman2020regression}
Gelman, Andrew, Jennifer Hill, and Aki Vehtari. 2020. \emph{Regression
and Other Stories}. Cambridge, UK: Cambridge University Press.

\bibitem[\citeproctext]{ref-mcelreath2020rethinking}
McElreath, Richard. 2020. \emph{Statistical Rethinking: A {Bayesian}
Course with Examples in {R} and {Stan}}. 2nd ed. Boca Raton, FL: CRC
Press.

\end{CSLReferences}


\backmatter


\end{document}
